\documentclass[titlepage, a4paper, openbib, 10pt]{article}

%#####################################
%Usepackages en installingen
\usepackage[top=1in, bottom=1in, left=1in, right=1in]{geometry}
\usepackage[pdftex]{graphicx}
\usepackage{fancyhdr}
\usepackage{sectionbox}
\usepackage[english]{babel}
\usepackage{chngcntr}
\usepackage{cite}
\usepackage{url}
\usepackage{makeidx}
\usepackage{paralist}
\usepackage{enumitem}
\usepackage{tocloft}
\usepackage{listliketab}	
\usepackage[table]{xcolor}
\usepackage{tabularx}
\usepackage{epsfig}
\usepackage{pdflscape}
\usepackage{pdfpages}
\usepackage{float}
\usepackage{multirow} 
\usepackage{rotating}
\usepackage[utf8]{inputenc}
\usepackage{color}
\usepackage{fp}
\usepackage[hidelinks]{hyperref}
\hypersetup{
    colorlinks=false,
    linkcolor=black,
    filecolor=black,
    urlcolor=black,
}
%\usepackage{draftwatermark}
%\SetWatermarkText{\textsc{Draft}}
%\SetWatermarkScale{5}
\newcommand{\red}[1]{
\textcolor{red}{#1}
}
\usepackage{listings}
\lstset{language=C,
basicstyle=\ttfamily\footnotesize,
frame=shadowbox,
escapeinside={(*@}{@*)},
mathescape=true,
showstringspaces=false,
showspaces=false,
breaklines=true}


%#####################################
%Glossary
%\usepackage{glossaries}
\usepackage{mfirstuc}
\usepackage[toc,nonumberlist,nopostdot]{glossaries}
\makeglossaries
\newglossaryentry{fpvsimp}{
      name=fpvsimp
    , description={The learning objective stating that at the end of this course: the student \textbf{understands} concepts of functional vs imperative semantics.}
    , first={\textbf{understands} the fundamental semantic difference between functional and imperative programming. \texttt{(FP VS IMP)}}
    , symbol=\texttt{FP VS IMP}
    }


\newglossaryentry{red}{
	name=red
	, description={The learning objective stating that at the end of this course: the student \textbf{understands} reduction strategies such as $\rightarrow_\beta$.}
	, first={\textbf{understands} reduction strategies such as $\rightarrow_\beta$. \texttt{(RED)}}
	, symbol=\texttt{RED}
}


\newglossaryentry{typ}{
	name=typ
	, description={The learning objective stating that at the end of this course: the student \textbf{understands} the basics of a functional type system.}
	, first={\texttt{understands} the basics of a functional type system. \texttt{(TYP)}}
, symbol=\texttt{TYP}
}


\newglossaryentry{fpext}{
	name=fpext
	, description={The learning objective stating that at the end of this course: the student \textbf{can program} with the common extensions of a functional programming language wrt the basic lambda calculus, such as \texttt{let}, \texttt{if}, \texttt{let-rec}, unions, tuples, records, etc. The languages of focus are F\# and Haskell.}
	, first={\texttt{can program} with the typical constructs of a modern functional language. The languages of focus are F\# and Haskell. \texttt{(FP EXT)}}
, symbol=\texttt{FP EXT}
}

%\loadglsentries{Glossary}


%\usepackage{showframe} %tmp
%#####################################
%Nieuwe commando's
\newcommand{\HRule}{\rule{\linewidth}{1pt}}
\newcommand{\organisatie}{\uppercase{Hogeschool Rotterdam / CMI}}
\newcommand{\modulenaam}{Functional programming fundamentals}
\newcommand{\modulecode}{\uppercase{INFSEN02-1}}
\newcommand{\studiejaar}{\uppercase{2015-2016}}
\newcommand{\stdPunten}{4}
\renewcommand{\author}{Tony Busker \& Giuseppe Maggiore}

\definecolor{lichtGrijs}{RGB}{169,169,169}



%#####################################
%Index en styling
\setlength{\cftbeforesecskip}{10pt}
\setlength\parindent{0pt}
\makeindex
\graphicspath{{../Img/}}
\counterwithin{figure}{subsection}
\pagestyle{fancy}
\setcounter{secnumdepth}{5}
\setcounter{tocdepth}{5}

%#####################################
%     Alles voor header/footer
\fancyhf[HL]{\nouppercase{\textit{\leftmark}}}
\setlength{\headheight}{36pt}
\lhead{\uppercase{\footnotesize Course description}}
\chead{\footnotesize \organisatie}
\rhead{\includegraphics[width=0.09\textwidth]{logo}}

\lfoot{\scriptsize \modulenaam}
\cfoot{\scriptsize \today}
\rfoot{\small \thepage}

\renewcommand{\headrulewidth}{0.4pt}
\renewcommand{\footrulewidth}{0.4pt}
%#####################################

\begin{document}


%#####################################
%Titlepage
\begin{titlepage}
\thispagestyle{fancy}
\ 
\vspace{5cm}

\begin{center}

	
	\Large \textbf \organisatie
	
	\vspace{1.5cm}
	
	\HRule \\[0.4cm]
	
	\Huge \textbf \modulenaam
	
	\vspace{1.7cm}
	
	\Large \textbf  \modulecode \\
	\studiejaar
	
	\vspace{0.4cm}
	
	\HRule \\[1.5cm]
\end{center}
\vfill

% Author and supervisor
\begin{tabular}{l l}
	Number of study points:  & \stdPunten{} ects\\
	Course owners: & \author\\
\end{tabular}


\end{titlepage}

%####### Contentpagina ########
%\renewcommand{\baselinestretch}{1.5}\normalsize
%\tableofcontents
%\newpage
%\listoffigures
%\newpage
%\listoftables
%\newpage

%########### Inhoud ###########

\shadowsectionbox
\section*{Module description}
\begin{tabularx}{\textwidth}{|>{\columncolor{lichtGrijs}} p{.26\textwidth}|X|}
	\hline
	\textbf{Module name:} & \modulenaam\\

	\hline
	\textbf{Module code: }& \modulecode\\
	\hline
	\textbf{Study points \newline and hours of effort:} & This module gives \stdPunten{}  ects, in correspondence with \FPeval{\result}{clip(\stdPunten*28)}\result{} hours:
	\begin{itemize}
		\item 2 X 3 x 6 hours of combined lecture and practical
		\item the rest is self-study
	\end{itemize} \\
	\hline
	\textbf{Examination:} & Written examination and practicums (with oral check) \\
	\hline
	\textbf{Course structure:} & Lectures, self-study, and practicums \\
	\hline
	\textbf{Prerequisite knowledge:} & all INFDEV courses. \\
	\hline
	\textbf{Learning materials:}  &
		\begin{itemize}
			\item Book: Friendly F\#; authors Giuseppe Maggiore and Giulia Costantini, available from the school library
			\item Book: Learn you a Haskell for a great good; available online for free
			\item Slides: found on N@tschool and on the GitHub repository \href{https://github.com/hogeschool/INFSEN01-1}{github.com/hogeschool/INFSEN01-1}
			\item exercises and assignments, to be done at home and during the practical part of the lectures (pdf): found on N@tschool and on the GitHub repository \href{https://github.com/hogeschool/INFDEV-Homework}{github.com/hogeschool/INFDEV-Homework}
		\end{itemize} \\
	\hline
	\textbf{Connected to competences:} & realiseren en ontwerpen \\
	\hline
	\textbf{Learning objectives:} &
		At the end of the course, the student:
			\begin{itemize}
                \item \glsfirst{fpvsimp}
                \item \glsfirst{red}
                \item \glsfirst{typ}
                \item \glsfirst{fpext}
			\end{itemize} \\
	\hline
%\end{tabularx}
%\newpage
%
%\begin{tabularx}{\textwidth}{|>{\columncolor{lichtGrijs}} p{.26\textwidth}|X|}
%	\hline
%	\textbf{Content:}&
%	\begin{itemize}
%		\item _
%	\end{itemize} \\
%	\hline
	\textbf{Course owners:} & \author\\
	\hline
	\textbf{Date:} & \today \\
	\hline
\end{tabularx}
%\newpage


\newpage
\section{General description}

Functional programming and functional programming languages are increasing in popularity for multiple reasons and in multiple ways, to the point that even mainstream languages such as Python, C++, C\#, and Java are being extended with more and more functional programming features such as tuples, lambda's, higher order functions, and even monads such as LINQ and async/await. Whole architectures such as the popular map/reduce are strongly inspired by functional programming.

``Java™ developers should learn functional paradigms now, even if they have no immediate plans to move to a functional language such as Scala or Clojure. Over time, all mainstream languages will become more functional'' [IBM].

``LISP is worth learning for a different reason — the profound enlightenment experience you will have when you finally get it. That experience will make you a better programmer for the rest of your days, even if you never actually use LISP itself a lot.'' – Eric S. Raymond

``SQL, Lisp, and Haskell are the only programming languages that I've seen where one spends more time thinking than typing.'' – Philip Greenspun

``I do not know if learning Haskell will get you a job. I know it will make you a better software developer.'' – Larry O’ Brien

The reason for this growth is to be found in the safe and deep expressive power of functional languages, which are capable of recombining simpler elements into powerful, complex other elements with less space for mistakes and more control in the hands of the programmer. This comes at a fundamental cost: functional languages are structurally different from imperative and object oriented languages, and thus a new mindset is required of the programmer that wishes to enter this new world. Moreover, functional languages often require more thought and planning, and are thus experienced, especially by beginners, as somewhat less flexible and supporting of experimentation.

\subsection{Relationship with other didactic units}
This module completes and perfects the understanding and knowledge of programming that was set up in the preceding INFDEV courses.

\section{Course program}
The course is structured into eight lectures.
The eight lectures take place during the eight weeks of the course, but are not necessarily in a one-to-one correspondance with the course weeks.

\subsection{Chapter 1 - foundations (weeks 1, and 2)}
\paragraph*{Topics}
\begin{itemize}
	\item Recap of imperative programming language semantics: the \textit{shared memory} model;
	\item Functional programming concepts: the lambda calculus and beta reduction;
	\item Adding static typing: the simply typed lambda calculus;
	\item Making the language usable: delta rules.
\end{itemize}

\subsection{Chapter 2 - practical applications (weeks 3, 4, 5)}
\paragraph*{Topics}			
\begin{itemize}
	\item From the lambda calculus to F\# (with practical lecture);
	\item Lazy evaluation;
	\item From the lambda calculus to Haskell (with practical lecture);
	\item Advanced constructs: lists (and list comprehensions), records, tuples, discriminated unions;
	\item Interop with other languages.
\end{itemize}

\subsection{Chapter 3 - patterns and practice (weeks 6, 7, 8)}
\paragraph*{Topics}			
\begin{itemize}
	\item Traversable entities that self-update;
	\item Monads introduction;
	\item Traversable entities that self-update with coroutines (with practical lecture).
\end{itemize}

\newpage
\section{Assessment}
The course is tested with two exams:
A series of assignments which have to be handed in, but will not be graded offline. There will be an oral/practicum check, which is based on the assignments,
and a written exam. The final grade is determined as follows: \\

\texttt{if exam grade $ >= 5.0 $ then practicum-grade else 0}

\paragraph*{Motivation for grade}
A professional software developer is required to be able to program code which is, at the very least, \textit{correct}.

In order to produce correct code, we expect students to show:
\begin{inparaenum}[\itshape i\upshape)]
\item a foundation of knowledge about how the semantics of the programming language actually work;
\item fluency when actually writing the code.
\end{inparaenum}

The quality of the programmer is ultimately determined by his actual code-writing skills, therefore the written exam will contain require you to write code. This ensures that each student is able to show that his work is his own and that he has adequate understanding of its mechanisms.



\subsection{Theoretical examination \modulecode}
The general shape of an exam for \texttt{\modulecode} is made up of a short series of highly structured open questions.
In each exam the content of the questions will change, but the structure of the questions will remain the same.
For the structure (and an example) of the theoretical exam, see the appendix.


\subsection{Practical examination \modulecode}
There is only one assignment, which is mandatory, and formatively assessed for feedback.

\begin{itemize}
  \item All assignments are to be uploaded to N@tschool or Classroom in the required space (Inlevermap or assignment);
  \item Each assignment is designed to assess the students knowledge related to one or more learning goal.
          If the teacher is unable te assess the students' ability related to the appropriate learning goal based on his work, then no points will be awarded for that part.
  \item \textit{The teachers still reserves the right to check the practicums handed in by each student, and to use it for further evaluation.}
  \item The university rules on fraude and plagiarism (Hogeschoolgids art. 11.10 -- 11.12) also apply to code;
  
The practical exam requires to complete one of the following assignments.

\begin{itemize}
	\item A 2D simulation of a supermarket with customers, cash registers, and various aisles
	\item A 2D simulation of a supply chain with trucks, containers, and ships
	\item An interpreter for a Python-like language (with a parser for an extra challenge)
	\item An interpreter for the lambda calculus  (with a parser for an extra challenge)
	
\end{itemize}

\end{itemize}

\newpage
\section*{Structure of exam \modulecode}
The general shape of a theoretical exam for \texttt{DEV 3} is made up of only two, highly structured open questions.

\paragraph{Question 1: } \ \\

\textbf{General shape of the question:} \textit{Given the following lambda program, and a series of delta rules, show the beta reductions for this program.}

\textbf{Concrete example of question:}

\textbf{Program:}
\lstset{numbers=left,basicstyle=\ttfamily\small}\lstset{language=[Sharp]C}\begin{lstlisting}
(((TRUE $\wedge$ TRUE) T) F)
\end{lstlisting}

\textbf{Delta rules:}
\lstset{numbers=left,basicstyle=\ttfamily\small}\lstset{language=[Sharp]C}\begin{lstlisting}
TRUE $\equiv$ $\lambda$t$\rightarrow$f$\rightarrow$t
FALSE $\equiv$ $\lambda$t$\rightarrow$f$\rightarrow$f
$\wedge$ $\equiv$ $\lambda$p$\rightarrow$q$\rightarrow$((p q) p)
\end{lstlisting}

\textbf{Concrete example of answer:}

\lstset{basicstyle=\ttfamily\small}\lstset{numbers=none}\lstset{language=ML}\begin{lstlisting}
(((TRUE $\wedge$ TRUE) T) F)

(((( (*@\underline{$\wedge$}@*) TRUE) TRUE) T) F)

((((($\lambda$p$\rightarrow$q$\rightarrow$((p q) p)) TRUE) TRUE) T) F)

((((($\lambda$p$\rightarrow$q$\rightarrow$((p q) p)) (*@\underline{TRUE}@*)) TRUE) T) F)

((((($\lambda$p$\rightarrow$q$\rightarrow$((p q) p)) ($\lambda$t$\rightarrow$f$\rightarrow$t)) TRUE) T) F)

((((($\lambda$p$\rightarrow$q$\rightarrow$((p q) p)) ($\lambda$t$\rightarrow$f$\rightarrow$t)) (*@\underline{TRUE}@*)) T) F)

((((($\lambda$p$\rightarrow$q$\rightarrow$((p q) p)) ($\lambda$t$\rightarrow$f$\rightarrow$t)) ($\lambda$t$\rightarrow$f$\rightarrow$t)) T) F)

((((*@\underline{(($\lambda$p$\rightarrow$q$\rightarrow$((p q) p)) ($\lambda$t$\rightarrow$f$\rightarrow$t))}@*) ($\lambda$t$\rightarrow$f$\rightarrow$t)) T) F)

(((($\lambda$q$\rightarrow$((($\lambda$t$\rightarrow$f$\rightarrow$t) q) ($\lambda$t$\rightarrow$f$\rightarrow$t))) ($\lambda$t$\rightarrow$f$\rightarrow$t)) T) F)

(((*@\underline{(($\lambda$q$\rightarrow$((($\lambda$t$\rightarrow$f$\rightarrow$t) q) ($\lambda$t$\rightarrow$f$\rightarrow$t))) ($\lambda$t$\rightarrow$f$\rightarrow$t))}@*) T) F)

((((($\lambda$t$\rightarrow$f$\rightarrow$t) ($\lambda$t$\rightarrow$f$\rightarrow$t)) ($\lambda$t$\rightarrow$f$\rightarrow$t)) T) F)

((((*@\underline{(($\lambda$t$\rightarrow$f$\rightarrow$t) ($\lambda$t$\rightarrow$f$\rightarrow$t))}@*) ($\lambda$t$\rightarrow$f$\rightarrow$t)) T) F)

(((($\lambda$f$\rightarrow$t$\rightarrow$f$\rightarrow$t) ($\lambda$t$\rightarrow$f$\rightarrow$t)) T) F)

(((*@\underline{(($\lambda$f$\rightarrow$t$\rightarrow$f$\rightarrow$t) ($\lambda$t$\rightarrow$f$\rightarrow$t))}@*) T) F)

((($\lambda$t$\rightarrow$f$\rightarrow$t) T) F)

((($\lambda$t$\rightarrow$f$\rightarrow$t) T) F)

((*@\underline{(($\lambda$t$\rightarrow$f$\rightarrow$t) T)}@*) F)

(($\lambda$f$\rightarrow$T) F)

(*@\underline{(($\lambda$f$\rightarrow$T) F)}@*)

T
\end{lstlisting}
	

\textbf{Points:} \textit{4 (50\% of total).}

\textbf{Grading:} \textit{Full points for more all correct steps and result. Half points if correct result is found with some reduction mistakes. Zero points otherwise.}

\textbf{Associated learning objective:} \glsfirst{red}

\ \\ 

\paragraph{Question 2: } \ \\

\textbf{General shape of question:} \textit{Given the following lambda calculus program, and a series of typing rules, give the full typing derivation for the program.}

\textbf{Concrete example of question:} 

\lstset{numbers=left,basicstyle=\ttfamily\small}\lstset{language=[Sharp]C}
\begin{lstlisting}
TODO
\end{lstlisting}

\textbf{Concrete example of answer:} \textit{}

\begin{lstlisting}
TODO
\end{lstlisting}

\textbf{Points:} \textit{4 (50\% of total).}

\textbf{Grading:} \textit{Full points for fully correct type derivation. Half points for minor mistakes but correct overall structure. Zero points otherwise.}

\textbf{Associated learning objective:} \glsfirst{typ}

\ \\

\newpage

\newpage
%\bibliographystyle{plain}
%\bibliography{references}
\glossarystyle{altlist}
\printglossaries
\newpage

\section*{Appendix 1: Assessment matrix}
	\begin{tabular}{|p{2cm}|p{4cm}|}
		\hline
		Learning objective & Dublin descriptors \\
		\hline
		\glssymbol{fpvsimp} & 1, 4, 5 \\
		\hline
        \glssymbol{red} & 1, 2, 4, 5 \\
        \hline
        \glssymbol{typ}& 1, 4, 5 \\
        \hline
        \glssymbol{fpext} & 1, 2 \\
        \hline
        \glssymbol{monads} & 1, 2 \\
        \hline
	\end{tabular}
	
	\vspace{1cm}

	Dublin-descriptors:
	\begin{enumerate}
		\item Knowledge and understanding
		\item Applying knowledge and understanding
		\item Making judgments
		\item Communication
		\item Learning skills
	\end{enumerate}


%\newpage
%\input{tex/Bijlage2}
%\newpage
%\input{tex/Bijlage3}
\printindex


\end{document}

